\documentclass{report}
\usepackage{amsmath,amssymb}
\setlength{\parindent}{0mm}
\setlength{\parskip}{1em}
\begin{document}
\begin{center}
\rule{6in}{1pt} \
{\large William Taitano \\
{\bf iFP: An optimal, fully implicit, fully conservative 1D2V Vlasov-Rosenbluth-Fokker-Planck code for ICF capsule implosion simulations}}

P O Box 1663 \\ MS-K717 \\ Los Alamos \\ NM \\ 87545
\\
{\tt taitano@lanl.gov}\\
Luis Chacon\\
Andrei Simakov\end{center}

Contrary to predictions of radiation-hydrodynamics design codes, the
National Ignition Facility has not succeeded in achieving ignition.
Recent experimental evidence suggests that plasma kinetic effects may
play an important role during Inertial Confinement Fusion (ICF) capsules�
implosion. Consequently, kinetic models and simulations may need to be
used to better understand experimental results and design ICF targets. We
present a new, optimal, fully implicit, and fully conservative 1D2V
Vlasov-Fokker-Planck (VFP) code, iFP, which simulates ICF implosions
kinetically. Such simulations are difficult to perform because of the
disparate time and length scales involved. The challenge in obtaining a
credible solution is further complicated by the need to enforce discrete
conservation properties.

In our studies, we employ the Rosenbluth formulation for the
Fokker-Planck collision operator. Our approach uses a fully implicit
temporal advance to step over stiff collision-time scales. For the
solver, we use a Jacobian-free Newton-Krylov method with an optimal
multigrid-based preconditioning technology. To address the issues of
velocity disparity between various species as well as those associated
with temporal and spatial temperature variations, we have developed: 1) a
velocity space meshing scheme, which adapts to the species� local thermal
velocity; and 2) an asymptotic expansion of the Rosenbluth potentials
based on the large ratio of thermal speeds of the fast-to-slow species
($v_{th,f}/v_{th,s} \gg 1$). We have also implemented a Lagrangian mesh,
which allows the physical space mesh to move as the capsule compresses.
Finally, we enforce discrete conservation of mass, momentum, and energy
by solving a set of discrete nonlinear constraints, which are derived
from continuum symmetries present in the VFP equations. Herein, we
present some results testing the code capabilities and show preliminary
simulations of a plasma shock.


\end{document}
